\documentclass[10pt,a4paper]{amsart}
\usepackage{amsmath,amssymb,amsthm,graphicx}
\usepackage{verbatim}
\newenvironment{reference}{\comment}{\endcomment}
%\newenvironment{reference}{}{}
\newenvironment{slogan}{\comment}{\endcomment}
\newenvironment{history}{\comment}{\endcomment}

% For commutative diagrams we use Xy-pic
\usepackage[all]{xy}

% We use 2cell for 2-commutative diagrams.
\xyoption{2cell}
% Theorem environments.
%
\theoremstyle{plain}
\newtheorem{theorem}[subsection]{Theorem}
\newtheorem{proposition}[subsection]{Proposition}
\newtheorem{lemma}[subsection]{Lemma}

\theoremstyle{definition}
\newtheorem{definition}[subsection]{Definition}
\newtheorem{example}[subsection]{Example}
\newtheorem{exercise}[subsection]{Exercise}
\newtheorem{situation}[subsection]{Situation}

\theoremstyle{remark}
\newtheorem{remark}[subsection]{Remark}
\newtheorem{remarks}[subsection]{Remarks}

\numberwithin{equation}{subsection}

% Macros
%
\def\lim{\mathop{\mathrm{lim}}\nolimits}
\def\colim{\mathop{\mathrm{colim}}\nolimits}
\def\Spec{\mathop{\mathrm{Spec}}}
\def\Hom{\mathop{\mathrm{Hom}}\nolimits}
\def\Ext{\mathop{\mathrm{Ext}}\nolimits}
\def\SheafHom{\mathop{\mathcal{H}\!\mathit{om}}\nolimits}
\def\SheafExt{\mathop{\mathcal{E}\!\mathit{xt}}\nolimits}
\def\Sch{\mathit{Sch}}
\def\Mor{\mathop{\mathrm{Mor}}\nolimits}
\def\Ob{\mathop{\mathrm{Ob}}\nolimits}
\def\Sh{\mathop{\mathit{Sh}}\nolimits}
\def\NL{\mathop{N\!L}\nolimits}
\def\CH{\mathop{\mathrm{CH}}\nolimits}
\def\proetale{{pro\text{-}\acute{e}tale}}
\def\etale{{\acute{e}tale}}
\def\QCoh{\mathit{QCoh}}
\def\Ker{\mathop{\mathrm{Ker}}}
\def\Im{\mathop{\mathrm{Im}}}
\def\Coker{\mathop{\mathrm{Coker}}}
\def\Coim{\mathop{\mathrm{Coim}}}

% Boxtimes
%
\DeclareMathSymbol{\boxtimes}{\mathbin}{AMSa}{"02}

%
% Macros for moduli stacks/spaces
%
\def\QCohstack{\mathcal{QC}\!\mathit{oh}}
\def\Cohstack{\mathcal{C}\!\mathit{oh}}
\def\Spacesstack{\mathcal{S}\!\mathit{paces}}
\def\Quotfunctor{\mathrm{Quot}}
\def\Hilbfunctor{\mathrm{Hilb}}
\def\Curvesstack{\mathcal{C}\!\mathit{urves}}
\def\Polarizedstack{\mathcal{P}\!\mathit{olarized}}
\def\Complexesstack{\mathcal{C}\!\mathit{omplexes}}
% \Pic is the operator that assigns to X its picard group, usage \Pic(X)
% \Picardstack_{X/B} denotes the Picard stack of X over B
% \Picardfunctor_{X/B} denotes the Picard functor of X over B
\def\Pic{\mathop{\mathrm{Pic}}\nolimits}
\def\Picardstack{\mathcal{P}\!\mathit{ic}}
\def\Picardfunctor{\mathrm{Pic}}
\def\Deformationcategory{\mathcal{D}\!\mathit{ef}}
\usepackage{cleveref}
\crefname{originallemma}{original lemma}{original lemmas}
\Crefname{originallemma}{Original lemma}{Original lemmas}
\makeatletter
\newlabel{lemma-faithful}{{04M6}{source}{}{sourcepointer-1}{}}
\newlabel{lemma-faithful@cref}{{[originallemma][0][]04M6}{[1][0][]source}{}{}{}}
\makeatother
\crefname{originalequation}{original equation}{original equations}
\Crefname{originalequation}{Original equation}{Original equations}
\makeatletter
\newlabel{equation-localize-at-scheme-ringed}{{04M1}{source}{}{sourcepointer-2}{}}
\newlabel{equation-localize-at-scheme-ringed@cref}{{[originalequation][0][]04M1}{[1][0][]source}{}{}{}}
\makeatother
\crefname{originallemma}{original lemma}{original lemmas}
\Crefname{originallemma}{Original lemma}{Original lemmas}
\makeatletter
\newlabel{lemma-relocalize-morphism-at-schemes}{{04M5}{source}{}{sourcepointer-3}{}}
\newlabel{lemma-relocalize-morphism-at-schemes@cref}{{[originallemma][0][]04M5}{[1][0][]source}{}{}{}}
\makeatother
\crefname{originalreference}{original reference}{original references}
\Crefname{originalreference}{Original reference}{Original references}
\makeatletter
\newlabel{sites-modules-lemma-relocalize-morphism-ringed-topoi}{{04J8}{source}{}{sourcepointer-4}{}}
\newlabel{sites-modules-lemma-relocalize-morphism-ringed-topoi@cref}{{[originalreference][0][]04J8}{[1][0][]source}{}{}{}}
\makeatother
\crefname{originalreference}{original reference}{original references}
\Crefname{originalreference}{Original reference}{Original references}
\makeatletter
\newlabel{etale-cohomology-theorem-fully-faithful}{{04I7}{source}{}{sourcepointer-5}{}}
\newlabel{etale-cohomology-theorem-fully-faithful@cref}{{[originalreference][0][]04I7}{[1][0][]source}{}{}{}}
\makeatother
\crefname{originalreference}{original reference}{original references}
\Crefname{originalreference}{Original reference}{Original references}
\makeatletter
\newlabel{spaces-lemma-space-presentation}{{0262}{source}{}{sourcepointer-6}{}}
\newlabel{spaces-lemma-space-presentation@cref}{{[originalreference][0][]0262}{[1][0][]source}{}{}{}}
\makeatother
\title{The locally ringed etale topos construction is fully faithful on algebraic spaces over the integers}
\author{The Stacks Project Authors}
\date{stacks-project:spaces-properties.tex:local-snapshot-20261012; GNU FDL 1.2 or later}
\begin{document}
\maketitle
\noindent\textit{Editorial note.} All delivered mathematical text is the author's own, copied verbatim from the pinned Stacks Project chapter file spaces-properties.tex; no statement or proof step was rewritten, reordered or removed.\par\medskip
\noindent\textit{Editorial note.} Span recomputation by label pairing (LANE\_BRIEF 10.2): in this chapter the author's proof environment follows the statement environment as a standalone block; the spec covers the complete environment pair (statement span, proof span) as recorded in the pool candidate.\par\medskip
\noindent\textit{Editorial note.} Occupation scan: combined claim\_key + statement\_anchor + generic-line scan over all registration files (active and withdrawn) found 13 S10-geom published cards for spaces-properties.tex (prop-occupied.json, 2026-10-14); those 13 labels are excluded from this batch by design. For the 77 free labels: 0 occupied. The registry mentions field (7557 files) and focused-seeds (645) both point to the same S10-geom cards (dual-source check), no other occupation.\par\medskip
\ifdefined\hypertarget\hypertarget{sourcepointer-1}{}\fi
\noindent\textit{Source reference.} Original lemma 04M6: sources/stacks-chow/tags.txt, source lines 6032--6032. Cross-chapter standard prerequisite cited by the author inside the counted span; it is used as a pointer and its own proof is not part of this item.\par\medskip
\ifdefined\hypertarget\hypertarget{sourcepointer-2}{}\fi
\noindent\textit{Source reference.} Original equation 04M1: sources/stacks-chow/tags.txt, source lines 6027--6027. Cross-chapter standard prerequisite cited by the author inside the counted span; it is used as a pointer and its own proof is not part of this item.\par\medskip
\ifdefined\hypertarget\hypertarget{sourcepointer-3}{}\fi
\noindent\textit{Source reference.} Original lemma 04M5: sources/stacks-chow/tags.txt, source lines 6031--6031. Cross-chapter standard prerequisite cited by the author inside the counted span; it is used as a pointer and its own proof is not part of this item.\par\medskip
\ifdefined\hypertarget\hypertarget{sourcepointer-4}{}\fi
\noindent\textit{Source reference.} Original reference 04J8: sources/stacks-chow/tags.txt, source lines 5929--5929. Cross-chapter standard prerequisite cited by the author inside the counted span; it is used as a pointer and its own proof is not part of this item.\par\medskip
\ifdefined\hypertarget\hypertarget{sourcepointer-5}{}\fi
\noindent\textit{Source reference.} Original reference 04I7: sources/stacks-chow/tags.txt, source lines 5893--5893. Cross-chapter standard prerequisite cited by the author inside the counted span; it is used as a pointer and its own proof is not part of this item.\par\medskip
\ifdefined\hypertarget\hypertarget{sourcepointer-6}{}\fi
\noindent\textit{Source reference.} Original reference 0262: sources/stacks-chow/tags.txt, source lines 2835--2835. Cross-chapter standard prerequisite cited by the author inside the counted span; it is used as a pointer and its own proof is not part of this item.\par\medskip
\begin{theorem}
\label{theorem-fully-faithful}
Let $X$, $Y$ be algebraic spaces over $\Spec(\mathbf{Z})$.
Let
$$
(g, g^\sharp) :
(\Sh(X_\etale), \mathcal{O}_X)
\longrightarrow
(\Sh(Y_\etale), \mathcal{O}_Y)
$$
be a morphism of locally ringed topoi. Then there exists a
unique morphism of algebraic spaces $f : X \to Y$ such that
$(g, g^\sharp)$ is isomorphic to $(f_{small}, f^\sharp)$.
In other words, the construction
$$
\textit{Spaces}/\Spec(\mathbf{Z})
\longrightarrow \textit{Locally ringed topoi},
\quad
X \longrightarrow (X_\etale, \mathcal{O}_X)
$$
is fully faithful (morphisms up to $2$-isomorphisms on the right hand side).
\end{theorem}

\begin{proof}
The uniqueness we have seen in
Lemma \ref{lemma-faithful}.
Thus it suffices to prove existence.
In this proof we will freely use the identifications of
Equation (\ref{equation-localize-at-scheme-ringed})
as well as the result of
Lemma \ref{lemma-relocalize-morphism-at-schemes}.

\medskip\noindent
Let $U \in \Ob(X_\etale)$, let
$V \in \Ob(Y_\etale)$
and let $s \in g^{-1}h_V(U)$ be a section. We may think of
$s$ as a map of sheaves $s : h_U \to g^{-1}h_V$. By
Modules on Sites,
Lemma \ref{sites-modules-lemma-relocalize-morphism-ringed-topoi}
we obtain a commutative diagram of morphisms of ringed topoi
$$
\xymatrix{
(\Sh(X_\etale/U), \mathcal{O}_U)
\ar[rr]_-{(j, j^\sharp)} \ar[d]_{(g_s, g_s^\sharp)} & &
(\Sh(X_\etale), \mathcal{O}_X) \ar[d]^{(g, g^\sharp)} \\
(\Sh(V_\etale), \mathcal{O}_V) \ar[rr] & &
(\Sh(Y_\etale), \mathcal{O}_Y).
}
$$
By
\'Etale Cohomology, Theorem \ref{etale-cohomology-theorem-fully-faithful}
we obtain a unique morphism of schemes $f_s : U \to V$ such that
$(g_s, g_s^\sharp)$ is $2$-isomorphic to $(f_{s, small}, f_s^\sharp)$.
The construction $(U, V, s) \leadsto f_s$ just explained satisfies
the following functoriality property: Suppose given morphisms
$a : U' \to U$ in $X_\etale$ and $b : V' \to V$ in $Y_\etale$
and a map $s' : h_{U'} \to g^{-1}h_{V'}$ such that the diagram
$$
\xymatrix{
h_{U'} \ar[d]_a \ar[r]_{s'} & g^{-1}h_{V'} \ar[d]^{g^{-1}b} \\
h_U \ar[r]^s & g^{-1}h_V
}
$$
commutes. Then the diagram
$$
\xymatrix{
U' \ar[r]_-{f_{s'}} \ar[d]_a & u(V') \ar[d]^{u(b)} \\
U \ar[r]^-{f_s} & u(V)
}
$$
of schemes commutes. The reason this is true is that the same condition
holds for the morphisms $(g_s, g_s^\sharp)$ constructed in
Modules on Sites,
Lemma \ref{sites-modules-lemma-relocalize-morphism-ringed-topoi}
and the uniqueness in
\'Etale Cohomology, Theorem \ref{etale-cohomology-theorem-fully-faithful}.

\medskip\noindent
The problem is to glue the morphisms $f_s$ to a morphism of algebraic
spaces. To do this first choose a scheme $V$ and a surjective \'etale
morphism $V \to Y$. This means that $h_V \to *$ is surjective and hence
$g^{-1}h_V \to *$ is surjective too. This means there exists a scheme $U$
and a surjective \'etale morphism $U \to X$ and a morphism
$s : h_U \to g^{-1}h_V$. Next, set $R = V \times_Y V$ and
$R' = U \times_X U$. Then we get
$g^{-1}h_R = g^{-1}h_V \times g^{-1}h_V$ as $g^{-1}$ is exact.
Thus $s$ induces a morphism $s \times s : h_{R'} \to g^{-1}h_R$.
Applying the constructions above we see that we get a
commutative diagram of morphisms of schemes
$$
\xymatrix{
R' \ar@<1ex>[d] \ar@<-1ex>[d] \ar[rr]_{f_{s \times s}} & &
R \ar@<1ex>[d] \ar@<-1ex>[d] \\
U \ar[rr]^{f_s} & &
V
}
$$
Since we have $X = U/R'$ and $Y = V/R$ (see
Spaces, Lemma \ref{spaces-lemma-space-presentation})
we conclude that this diagram
defines a morphism of algebraic spaces $f : X \to Y$ fitting
into an obvious commutative diagram.
Now we still have to show that $(f_{small}, f^\sharp)$ is
$2$-isomorphic to $(g, g^\sharp)$.
Let $t_V : f_{s, small}^{-1} \to g_s^{-1}$ and
$t_R : f_{s \times s, small}^{-1} \to g_{s \times s}^{-1}$ be
the $2$-isomorphisms which are given to us by the construction above.
Let $\mathcal{G}$ be a sheaf on $Y_\etale$. Then we see that
$t_V$ defines an isomorphism
$$
f_{small}^{-1}\mathcal{G}|_{U_\etale}
=
f_{s, small}^{-1}\mathcal{G}|_{V_\etale}
\xrightarrow{t_V}
g_s^{-1}\mathcal{G}|_{V_\etale}
=
g^{-1}\mathcal{G}|_{U_\etale}.
$$
Moreover, this isomorphism pulled back to $R'$ via either projection
$R' \to U$ is the isomorphism
$$
f_{small}^{-1}\mathcal{G}|_{R'_\etale}
=
f_{s \times s, small}^{-1}\mathcal{G}|_{R_\etale}
\xrightarrow{t_R}
g_{s \times s}^{-1}\mathcal{G}|_{R_\etale}
=
g^{-1}\mathcal{G}|_{R'_\etale}.
$$
Since $\{U \to X\}$ is a covering in the site $X_{spaces, \etale}$
this means the first displayed isomorphism descends to an isomorphism
$t : f_{small}^{-1}\mathcal{G} \to g^{-1}\mathcal{G}$
of sheaves (small detail omitted). The isomorphism is functorial
in $\mathcal{G}$ since $t_V$ and $t_R$ are transformations of functors.
Finally, $t$ is compatible with $f^\sharp$ and $g^\sharp$ as
$t_V$ and $t_R$ are (some details omitted).
This finishes the proof of the theorem.
\end{proof}

\end{document}
